\chapter{Đánh giá và nhìn lại}
\label{ch:danh-gia}

\section{So sánh với mục tiêu ban đầu}

\begin{table}[H]
  \centering
  \small
  \begin{tabularx}{\textwidth}{X l X}
    \toprule
    \textbf{Mục tiêu (Lời mở đầu)} & \textbf{Kết quả} & \textbf{Minh chứng} \\
    \midrule
    6 mục lõi Tầng 1A & Đạt 6/6 & Bảng~\ref{tab:tong-hop-1a} \\
    Custom node parser, persistent index, trích nguồn & Đạt & Chương~\ref{ch:node-parser}--\ref{ch:trich-dan} \\
    Từ chối câu ngoài corpus & Một phần & Chưa đo tỷ lệ đúng \\
    Thí nghiệm chunking & Chưa & Mục~\ref{sec:node} \\
    Hybrid, rerank (Tầng 3) & Đạt, có số liệu & Chương~\ref{ch:nang-cao} \\
    Kiểm thử + CI (Tầng 2) & Đạt & 26 kiểm thử \\
    Chấm câu trả lời tự động & Chưa chạy & Mục~\ref{sec:danh-gia-tra-loi} \\
    \bottomrule
  \end{tabularx}
  \caption{Đối chiếu mục tiêu và kết quả.}
\end{table}

\paragraph{Các quyết định đổi hướng.}
\begin{itemize}
  \item \textbf{Từ OCR theo trang sang Markdown hai tầng}: vì Điều bị chia đôi giữa hai trang và Điều của
        quyết định ban hành trùng nhãn (Chương~\ref{ch:node-parser}).
  \item \textbf{Từ một văn bản sang sáu văn bản đang áp dụng}: để trả lời được nội dung sửa đổi; văn bản hết
        hiệu lực bị loại vì lấn át Quy chế 2021 khi truy hồi.
  \item \textbf{Chọn profile \texttt{server} làm cấu hình chính, LLM chạy từ xa}: GPU 4\,GB không chạy nổi
        \texttt{qwen3:8b} cùng embedding; LLM chạy trên máy GPU thuê hoặc Colab
        (\texttt{scripts/setup\_server\_v2.sh}), embedding chạy local ở fp16 hoặc CPU.
  \item \textbf{Tắt rerank và hybrid}: đo thấy không giúp ở quy mô 54 Node (Chương~\ref{ch:nang-cao}).
\end{itemize}

\section{Hạn chế}

\begin{itemize}
  \item \textbf{Độ mạnh thống kê của bộ câu hỏi.} 31 + 10 = 41 câu là ít; khoảng tin cậy rộng (bộ 02:
        R@1 90\% với CI 70--100\%) và kiểm định vector--BM25 chưa đạt $p<0{,}05$. Cần mở rộng lên
        khoảng 150--300 câu, phân tầng theo loại câu hỏi, và có hai người gán nhãn độc lập để báo cáo
        độ đồng thuận.
  \item \textbf{Corpus nhỏ.} 6 văn bản / 54 Node. Kết luận ``rerank và hybrid không giúp'' chỉ đúng ở
        quy mô này; ở kho vài trăm Node trở lên, hai kỹ thuật đó cần được đo lại.
  \item \textbf{Chất lượng OCR.} Văn bản quan trọng nhất (QĐ 3183/2025) có chất lượng OCR thấp nhất
        (94,9\% âm tiết hợp lệ) và vẫn còn lỗi như ``Sửa đồi''.
  \item \textbf{Metadata hiệu lực còn trống.} \texttt{\_metadata\_van\_ban.csv} chưa có số hiệu, ngày
        ban hành, tình trạng hiệu lực, nên chưa áp dụng được lọc theo hiệu lực hay theo khóa.
  \item \textbf{Đánh giá câu trả lời chưa chạy tự động.} Script đã có và LLM đã chạy được trên Colab,
        nhưng \texttt{make eval-answers} chưa chạy; câu trả lời mới chỉ được chấm tay.
  \item \textbf{Hạ tầng đo.} Số liệu trong chương này đo trên CPU; thời gian tuyệt đối không so sánh
        được với các con số trên GPU ở bản báo cáo trước.
  \item \textbf{Tầng sinh.} Không nhớ hội thoại, không có công cụ tra bảng/quy đổi điểm, chỉ một lượt
        gọi LLM trên ngữ cảnh đã truy hồi.
  \item \textbf{Không phân biệt khóa tuyển sinh.} Thang điểm của QĐ 3183 chỉ áp dụng từ khóa 2025, nhưng
        hệ thống không hỏi hay nêu khóa nên câu trả lời về điểm chỉ đúng với một nhóm sinh viên
        (Mục~\ref{sec:chay-lai-56-cau}).
  \item \textbf{Chú giải thang điểm trong prompt chưa được đo.} Đây là kiến thức riêng của quy chế SGU
        đặt trong prompt (Mục~\ref{sec:nham-thang-diem}), và chưa có phiên chạy nào xác nhận nó sửa được lỗi.
\end{itemize}


\section{Phân tích lỗi}

\subsection{Lỗi 1: LLM nhầm thang điểm hệ 10 với hệ 4}

\begin{itemize}
  \item \textbf{Hiện tượng}: ``A thì được bao nhiêu điểm hệ 10'' $\rightarrow$ ``A được quy đổi thành 4,0 điểm hệ 10''
        (nhật ký Gradio, ngày 08/10, hai lần hỏi).
  \item \textbf{Phát hiện}: đọc nhật ký \texttt{questions\_log.jsonl}; top-$k$ (Bảng~\ref{tab:top-k}) cho thấy truy
        hồi đúng, lỗi nằm ở tầng sinh.
  \item \textbf{Nguyên nhân và cách sửa}: Mục~\ref{sec:nham-thang-diem}. Chưa đo lại sau khi sửa.
\end{itemize}

\subsection{Lỗi 2: \texttt{CUDA out of memory} khi nạp embedding cùng reranker}

\begin{itemize}
  \item \textbf{Hiện tượng}: trên card 4\,GB, profile \texttt{server} chiếm khoảng 3,5\,GB sau khi nạp embedding
        và báo hết bộ nhớ khi nạp reranker.
  \item \textbf{Nguyên nhân}: \texttt{SentenceTransformerRerank} không nhận \texttt{dtype}, luôn nạp fp32.
  \item \textbf{Cách sửa}: nạp reranker trên CPU, ép fp16 rồi mới chuyển lên GPU:
\end{itemize}

\begin{lstlisting}[language=Python, caption={Nạp reranker fp16 (\texttt{make\_local\_reranker()} trong \texttt{src/retriever.py}).}, label={lst:rerank-fp16}]
reranker = SentenceTransformerRerank(model=cfg.reranker_model_name, top_n=cfg.reranker_top_n,
                                     device="cpu" if fp16 else device)
if fp16:
    reranker._model.half().to(device)
    reranker.device = device
\end{lstlisting}

Sau khi sửa, cả hai mô hình cùng nằm trên GPU, tổng khoảng 2,2\,GB (Mục~\ref{sec:phan-cung}).

\section{AI Disclosure}

\chualam{phần này viết ở bản báo cáo cuối, dựa trên nhật ký dùng AI của nhóm}. Nội dung sẽ gồm:
\begin{enumerate}
  \item Công cụ AI đã dùng và dùng vào việc gì.
  \item Phần có tỷ trọng AI sinh cao (ghi file / chức năng).
  \item Ít nhất hai tình huống AI trả kết quả \emph{sai}: prompt, AI trả về gì, sai ở đâu, phát hiện thế nào,
        sửa ra sao.
  \item Tự đánh giá: nếu bỏ AI, phần nào nhóm vẫn tự viết lại được.
\end{enumerate}
